% =====================================================================
%  thmsmith Stage -1 proposal skeleton.
%  Written automatically by the thmsmith TS pipeline before Stage -1 runs.
%  Locked parameters: qid=pid\_reversible\_dynamic\_spillover, specialization=v1
%
%  HOW TO USE THIS FILE
%  --------------------
%  - The preamble (everything above the `% --- BEGIN CONTENT ---` marker)
%    and the `\section{...}` headers below are fixed by the pipeline.
%    Do NOT rewrite or rename them; the Step 3 reviewer subagent and the
%    downstream Stage 0 / Stage 1 prompts grep for these section titles.
%  - Each section contains exactly one `% TODO(stage-1 ...): ...`
%    placeholder. Replace each TODO with the content described for that
%    section in `CausalSmith/tools/src/discovery/prompts/stage_neg1_2_draft.txt`
%    ("REQUIRED OUTPUT FORMAT (Stage -1 proposal .tex)").
%  - The `\paragraph{Locked parameters.}` block in the metadata block
%    must pin the chosen `(qid, specialization, p, assignment, target,
%    regression)` precisely. If the proposal maps to the Q1 pool-mode,
%    reproduce that block verbatim from
%    `CausalSmith/doc/panel_formula_question_generator_note_with_common_prompts.tex`
%    (Q2..Q6 templates have been retired). Otherwise (free-form
%    `(qid, specialization)` — the default), define each parameter
%    explicitly here (no implicit conventions). Stage 0 quotes this block;
%    do not rephrase it after locking.
%  - Removing a section header, renaming a macro, or rewriting the
%    preamble is a regression: the file must still match this skeleton
%    after you finish.
% =====================================================================

\documentclass[11pt]{article}

\usepackage{amsmath,amssymb,amsthm,mathtools}
\usepackage{enumitem}
\usepackage[colorlinks=true,linkcolor=blue,citecolor=blue,urlcolor=blue]{hyperref}
\usepackage{natbib}

\newcommand{\E}{\mathbb{E}}
\renewcommand{\Pr}{\mathbb{P}}
\newcommand{\one}{\mathbf{1}}
\newcommand{\transpose}{\mathsf{T}}
\newcommand{\calH}{\mathcal{H}}
\newcommand{\calF}{\mathcal{F}}
\newcommand{\calR}{\mathcal{R}}
\newcommand{\R}{\mathbb{R}}
\newcommand{\plim}{\operatorname*{plim}}

\theoremstyle{definition}
\newtheorem{definition}{Definition}
\newtheorem{assumption}{Assumption}
\newtheorem{example}{Example}

\theoremstyle{plain}
\newtheorem{theorem}{Theorem}
\newtheorem{conjecture}{Conjecture}
\newtheorem{proposition}{Proposition}
\newtheorem{lemma}{Lemma}
\newtheorem{corollary}{Corollary}

\theoremstyle{remark}
\newtheorem{remark}{Remark}

\title{thmsmith Stage -1 proposal: \texttt{pid\_reversible\_dynamic\_spillover} / \texttt{v1}%
  \\ \large\textit{A finite-window frontier for reversible-spell long-run partial-ID intervals}}
\author{thmsmith Stage -1 proposer}
\date{\today}

\begin{document}
\maketitle

% --- BEGIN CONTENT ---  (everything above is fixed by the pipeline)

\paragraph{Locked parameters.}
This is a partial-identification anchor, not a Q1 panel pool-mode.  The locked tuple is
\[
\begin{gathered}
(\texttt{qid},\texttt{specialization},P,\theta,
  \text{identifying restrictions},\text{bounding functional})\\
=
(\texttt{pid\_reversible\_dynamic\_spillover},\texttt{v1},
  P_{\infty},\theta_{\infty}(\rho),
  \mathcal R_{\mathrm{rev},m,\rho,\mathrm{Cesaro}},
  [\underline\theta_{\infty}(P_{\infty},\rho),
    \overline\theta_{\infty}(P_{\infty},\rho)]) .
\end{gathered}
\]
Here \(P_{\infty}\) is a projectively consistent law for an infinite discrete-time panel \((Y_t,D_t,X_t)_{t\ge 0}\), with \(D_t\in\{0,1\}\), finite \(X_t\), bounded outcomes \(Y_t\in[0,1]\), and reversible treatment paths.  The target
\[
\theta_{\infty}(\rho)=
\lim_{T\to\infty}\frac1T\sum_{t=1}^{T}
\E\{Y_t(\mathbf 1_m)-Y_t(\mathbf 0_m)\}
\]
is the long-run average effect of sustained treatment versus sustained no treatment when potential outcomes have \(m\)-period treatment memory and adjacent-spell carryover is bounded by the scalar \(\rho\).  The restriction class \(\mathcal R_{\mathrm{rev},m,\rho,\mathrm{Cesaro}}\) consists of consistency, bounded outcomes, finite \(m\)-memory, projective consistency of the observable laws, positive recurrent reversible support, identified counterfactual transition kernels for the two sustained regimes, and a global within-current-treatment spillover bound on latent reward maps.  The bounding functional is the Hausdorff/Cesaro limit, when it exists, of finite-horizon sharp intervals; the proposal asks when this limit equals the stationary Markov-reward support interval and when the interchange of horizon limits with sharp extrema fails.

\begin{abstract}
Finite-horizon dynamic treatment formulas do not by themselves define a long-run estimand when treatment can turn on, turn off, and turn on again.  This proposal asks for a finite-window frontier deciding when sharp partial-identification intervals for reversible bounded-carryover treatments converge to a stationary or Cesaro-limit identified set.  The routine background component is a Markov-reward support formula for long-run bounds once a stationary reward correspondence and sustained-regime transition kernels are already fixed.  The positive flagship conjecture is an if-and-only-if criterion: on a nontrivial class of projectively consistent reversible-spell laws, vanishing finite-window endpoint diameter is exactly equivalent to convergence of the finite sharp intervals to the stationary long-run interval, with stationary liftability following from the class restrictions.  The negative flagship conjecture is a block-switching sharpness witness showing that finite-horizon sharp-bounds workflows cannot be summarized by a single stationary interval unless that frontier criterion is reported.
\end{abstract}

\section{Introduction}
\paragraph{Why the problem is important.}
Many empirical treatments are reversible: medication use, training participation, sanctions, benefit eligibility, and local policy exposure all contain repeated starts and stops.  Finite-horizon dynamic treatment results identify or bound the effect of a specified treatment sequence, but an applied researcher often wants the long-run average effect of sustained treatment.  With bounded spillovers across adjacent spells, that long-run quantity is not just a large-\(T\) version of a two-period formula: endpoint faces of the finite-horizon identified set may move with the occupation distribution of treatment histories and with the active latent reward face.

\paragraph{The main finding (proposed).}
Proposition~\ref{thm:stationary-reward} records the routine stationary Markov-reward bound for a fixed long-run reward correspondence.  Conjecture~\ref{conj:cesaro-commutation} is the positive flagship proposal: for a class of projectively consistent reversible-spell laws with tight finite compatibility projections, the computable finite-window endpoint diameter is zero in the tail if and only if the finite sharp intervals converge to the stationary reward interval.  Conjecture~\ref{conj:face-cycling} is the diagnostic-necessity statement: a finite-state reversible panel with increasing blocks of latent carryover compatibility can make sharp endpoint subsequences separate even though \(Q_0\) and \(Q_1\) are uniformly geometrically ergodic.

\paragraph{What is new.}
\citet{Han2021} identifies finite dynamic treatment-sequence effects under period-by-period exclusion restrictions and emphasizes timing and state dependence, but does not give conditions under which finite-horizon identified sets define a long-run reversible-spell ATE.  \citet{ChernozhukovFernandezValHahnNewey2013} and \citet{JunLeeShin2016} prove sharp panel bounds under time-homogeneity or stationarity restrictions; their epsilon-gap is that stationarity is imposed as the identifying restriction, whereas this proposal asks for the exact finite-window frontier at which a sequence of reversible-spell finite sharpness problems becomes a stationary long-run bound.  \citet{Schneider1993} and \citet{RockafellarWets1998} cover generic compact-convex support-function and variational stability facts, so the causal novelty is not support-function continuity itself; it is the sharp latent-law question of whether the observed finite-horizon causal compatibility sets have a vanishing endpoint diameter and a liftable stationary face.  \citet{KalouptsidiKitamuraLimaSouzaRodrigues2020} and \citet{DuarteFinkelsteinKnoxMummoloShpitser2024} represent finite computational sharp-bounds workflows, while \citet{HanYang2024} is a computational MTE/policy-evaluation partial-identification comparator; the proposed gap is an analytic no-free-stationary-summary diagnostic for reversible spell histories, not another finite constrained optimization.

\paragraph{How it positions in the literature.}
The anchor cluster is partial identification with dynamic treatment timing and panel structure.  The proposal consumes the Robins g-formula only as the zero-uncertainty endpoint and treats assumption-free sustained-intervention bounds as the unrestricted endpoint.  Its core is closer to dynamic panel partial identification and Markov-reward average-cost theory, but it does not follow from the time-homogeneity shrinkage result in \citet{ChernozhukovFernandezValHahnNewey2013}: their identifying restriction directly links the latent structural objects across time, while here the same observed reversible panel may have finite sharpness faces that remain compatible at every horizon yet fail the endpoint Cauchy frontier.

\paragraph{Implications for the community.}
The concrete consumer is long-run policy evaluation for reversible treatment-history designs of the kind studied by \citet{DeChaisemartinDHaultfoeuille2024}: instead of reporting only finite lagged effects under parallel trends, an analyst could state whether a bounded-carryover sensitivity analysis has a valid long-run sharp interval at all.  The same result would give the missing convergence target for \citet{DiemerShiSwanson2024}-style sustained dynamic bounds, which currently emphasize unrestricted finite-horizon width rather than the stationary limit.

\paragraph{How research benefits from the conclusion.}
Researchers should report the finite-horizon-to-long-run stability check before interpreting reversible-spell dynamic bounds as a long-run average treatment effect.

\section{Related work}
\citet{Robins1986} introduces the g-formula for sustained time-varying treatments under sequential exchangeability; it is the zero-spillover, point-identification endpoint that any bounded-carryover long-run formula must respect.  \citet{Han2021} develops nonparametric identification for finite dynamic treatment sequences and is the closest finite-horizon dynamic-treatment comparator, but it does not characterize Cesaro limits of sharp reversible-spell intervals.  \citet{DiemerShiSwanson2024} gives assumption-free partial-identification bounds for sustained static and dynamic interventions; it supplies the unrestricted finite-horizon endpoint that the proposed long-run interval must nest.

\citet{ChernozhukovFernandezValHahnNewey2013} derive sharp average and quantile effect bounds in nonseparable panel models, including dynamic state-dependence settings under time-homogeneity; this proposal differs by using bounded spillover plus horizon-limit stability rather than time homogeneity as the source of identifying power.  \citet{JunLeeShin2016} give sharp identifiable bounds for potential-outcome distributions in panel data with dynamic treatment decisions under stationarity; the present proposal asks when stationarity of the limiting support functional follows from a finite-horizon problem rather than being imposed directly.  \citet{DeChaisemartinDHaultfoeuille2024} study non-binary, non-absorbing treatments with lags under a parallel-trends DID design; they are the applied reversible-lag consumer, while the proposed result is a non-DID partial-identification target.

\citet{HanYang2024} give a computational approach to partial identification of MTE-based policy evaluation with unobserved heterogeneity; it is relevant only as a finite-optimization partial-ID comparator, not as a dynamic-treatment or Markov-counterfactual bounds result.  \citet{DuarteFinkelsteinKnoxMummoloShpitser2024} automate sharp causal bounds in finite discrete settings by searching over admissible data-generating processes; this proposal is narrower but asks for the symbolic horizon-limit and finite-window endpoint-diameter diagnostics that a sequence of finite black-box programs would not state.  \citet{KalouptsidiKitamuraLimaSouzaRodrigues2020} compute sharp bounds for structural dynamic discrete-choice counterfactuals; their constrained optimization motivates the dynamic-counterfactual setting but relies on structural dynamic-choice primitives rather than a nonstructural bounded-spillover panel law, and their finite-horizon workflow is the named target for the no-free-stationary-summary warning in Conjecture~\ref{conj:face-cycling}.  \citet{HeckmanHumphriesVeramendi2016} decompose dynamic treatment effects into direct effects and continuation values; the proposal's stationary reward formula makes the continuation component explicit in a partial-identification rather than IV point-estimation setting.

\citet{Puterman1994} and \citet{Altman1999} give the standard average-reward and occupation-measure linear-programming toolkit behind the stationary support formula.  Their results explain why Proposition~\ref{thm:stationary-reward} is background once \(\mathcal G(P_{\infty},\rho)\) is fixed, but they do not address the causal sharpness problem that finite latent-law compatibility polytopes must project to the same stationary exposed face before the horizon limit can commute with endpoint optimization.  \citet{Schneider1993} and \citet{RockafellarWets1998} supply the closest mathematical support-function and Hausdorff-convergence tools; the proposed epsilon-gap is that those theorems do not decide whether the projected sets arise sharply from one projective causal panel law or whether reversible-spell carryover constraints can force horizon-indexed endpoint cycling.

In-catalogue CausalSmith prior art does not already contain this kernel.  The \texttt{doc/API.md} ExactID/PartialID catalogue records static Manski bounded-outcome wrappers under \texttt{ManskiNonparametricBounds\_Study}; those entries concern cross-sectional bounded response, not reversible dynamic histories or long-run limits.  The banked \texttt{pid\_dynamic\_iv\_compliance} proposals study retained instrument memory and endpoint-face witnesses for dynamic IV bounds, but not a projectively consistent reversible-spell Cesaro limit.  The banked \texttt{pid\_long\_term\_surrogate} and \texttt{pid\_did\_anticipation\_bounded} proposals use bounded transition or shared-coordinate templates, but their reviewer gaps concern finite witnesses or DID support functions rather than the limit/interchange criterion proposed here.

\section{Formal setup}
Let \(P_{\infty}\) be the law of an infinite panel process \((Y_t,D_t,X_t)_{t\ge0}\), where \(D_t\in\{0,1\}\), \(X_t\in\mathcal X\) for finite \(\mathcal X\), and \(Y_t\in[0,1]\).  For a memory length \(m\), write \(H_t=(D_{t-m+1},\ldots,D_t,X_t)\in\calH_m\), using the obvious initialization convention for \(t<m\).  The two sustained treatment histories are denoted \(\mathbf 1_m\) and \(\mathbf 0_m\), with state components \((\mathbf 1_m,x)\) and \((\mathbf 0_m,x)\).

Potential outcomes have finite memory: \(Y_t(\bar d)\) depends on \(\bar d\) only through the last \(m\) treatment statuses and the current state \(X_t\).  For \(a\in\{0,1\}\), let \(Q_a\) be the identified transition kernel of \(X_{t+1}\) under the sustained action \(D_t=a\), and let \(\pi_a\) be its stationary distribution when it exists.  A latent reward map is a pair \(r=(r_1,r_0)\), where \(r_a(x)\in[0,1]\) is the mean outcome at state \(x\) under the sustained \(a\)-history.  Observed reversible histories provide anchor intervals \(A_a(x;\rho)=[\ell_a(x;\rho),u_a(x;\rho)]\) for each \(r_a(x)\); these intervals are induced by consistency, finite memory, and the global adjacent-spell carryover bound \(\rho\).  The stationary reward correspondence is
\[
\mathcal G(P_{\infty},\rho)
=\{r:\ell_a(x;\rho)\le r_a(x)\le u_a(x;\rho)
  \text{ for all }a,x,\text{ and }r\text{ is dynamically compatible with }P_{\infty}\}.
\]
The long-run identified set is
\[
\Theta_{\infty}(P_{\infty},\rho)
=\left\{\sum_x\pi_1(x)r_1(x)-\sum_x\pi_0(x)r_0(x):
  r\in\mathcal G(P_{\infty},\rho)\right\}.
\]
For finite \(T\), \(\mathcal C_T(P_T,\rho)\) denotes the finite-horizon latent-law compatibility set: all laws for the relevant potential-outcome histories whose observed marginal is \(P_T\) and whose reward coordinates satisfy the \(\rho\)-carryover inequalities.  The sharp finite-horizon identified interval for the \(T\)-period average effect is
\(\Theta_T(P_T,\rho)=[\underline\theta_T,\overline\theta_T]\), the range of that average over \(\mathcal C_T(P_T,\rho)\).  Let \(\Pi_T:\mathcal C_T(P_T,\rho)\to\mathcal Z_T\) be the linear projection that records only occupation-reward coordinates
\[
z_T=(\nu_{T,1},\nu_{T,0},\bar r_{T,1},\bar r_{T,0}),
\]
where \(\nu_{T,a}\) is the empirical occupation measure of states while the current sustained arm is \(a\), and \(\bar r_{T,a}(x)\) is the projected mean reward coordinate for state \(x\) in arm \(a\).  The endpoint functional is the linear map
\[
L_T(z_T)=\sum_x\nu_{T,1}(x)\bar r_{T,1}(x)-\sum_x\nu_{T,0}(x)\bar r_{T,0}(x).
\]
All finite projected sets are embedded in the common compact metric space
\[
K=\Delta(\mathcal X)^2\times[0,1]^{2|\mathcal X|}
\]
with the product Euclidean metric.  Let \(\widehat{\mathcal Z}_T\subset K\) be the normalized image of \(\Pi_T\mathcal C_T(P_T,\rho)\), so that the occupation components are probability measures and \(L(z)=\nu_1r_1-\nu_0r_0\) is a common normalized objective on \(K\).  For \(s\in\{-,+\}\), define \(c^+=L\), \(c^-=-L\), the support value \(h_T^s=\sup_{z\in\widehat{\mathcal Z}_T}c^s(z)\), and the exposed endpoint face
\[
F_T^s=\arg\max_{z\in\widehat{\mathcal Z}_T}c^s(z).
\]
The finite-window endpoint diameter is
\[
\Delta_{T,J}
=\max_{T\le T',T''\le T+J}
\max_{s\in\{-,+\}}\left|h_{T'}^s-h_{T''}^s\right|,
\]
and its tail version is \(\Delta_T=\sup_{J\ge1}\Delta_{T,J}\).  This is the one-dimensional support criterion that can fail even when every finite horizon has a sharp nonempty interval.
For \(\varepsilon>0\), the observable exposed-face gap is
\[
\Gamma_T^s(\varepsilon)=
\inf\{h_T^s-c^s(z):z\in\widehat{\mathcal Z}_T,\ d(z,F_T^s)\ge\varepsilon\},
\]
with \(\Gamma_T^s(\varepsilon)=+\infty\) if the set is empty.  For a finite window \(J\ge1\), the computable stability diagnostic is
\[
\mathfrak D_{T,J}(\varepsilon)
=\max_{T\le T',T''\le T+J}
\left[
d_H(\widehat{\mathcal Z}_{T'},\widehat{\mathcal Z}_{T''})
+\max_{s\in\{-,+\}}d_H(F_{T'}^s,F_{T''}^s)
+\max_{s\in\{-,+\}}\left|\Gamma_{T'}^s(\varepsilon)-\Gamma_{T''}^s(\varepsilon)\right|
\right].
\]
The asymptotic tail diagnostic is \(\mathfrak D_T(\varepsilon)=\sup_{J\ge1}\mathfrak D_{T,J}(\varepsilon)\).  Each finite-window statistic \(\Delta_{T,J}\), \(\mathfrak D_{T,J}\), finite \(\widehat{\mathcal Z}_T\), face \(F_T^s\), support value \(h_T^s\), and gap \(\Gamma_T^s\) is determined by the finite observed law \(P_{T+J}\), the memory length \(m\), and the scalar bound \(\rho\).  The supremum over all future windows is a population tail property of \(P_{\infty}\), not a finite-window observable; empirical use would report a sequence of increasing-window checks \(\Delta_{T,J}\) and \(\mathfrak D_{T,J}\).
The stationary projection set is
\[
\mathcal Z_{\infty}
=\{(\pi_1,\pi_0,r_1,r_0):r\in\mathcal G(P_{\infty},\rho)\},
\]
with \(L(z)=\pi_1r_1-\pi_0r_0\) on \(\mathcal Z_\infty\).  Let \(F_\infty^s=\arg\max_{z\in\mathcal Z_\infty}c^s(z)\).  The stationary lift defect is
\[
\Lambda_T
=d_H(\widehat{\mathcal Z}_T,\mathcal Z_{\infty})
  +\max_{s\in\{-,+\}}d_H(F_T^s,F_\infty^s),
\]
with \(\Lambda_T=+\infty\) if an exposed stationary face is empty.  The proposed frontier criterion is
\[
\Psi_T=\Delta_T .
\]
The term \(\Lambda_T\) records stationary liftability but is not the finite-window test; Conjecture~\ref{conj:cesaro-commutation} states the class where the observable endpoint Cauchy criterion \(\Psi_T=\Delta_T\) is exactly the convergence frontier and where \(\Lambda_T\to0\) follows as a consequence.  If \(z_T\in\widehat{\mathcal Z}_T\) and \(z_\infty\in\mathcal Z_\infty\) share the same reward coordinates, their finite-horizon boundary correction is
\[
B_T(z_T,z_\infty)=L(z_T)-L(z_\infty),
\]
the endpoint change caused only by the finite-horizon occupation measure rather than by the reward face.  When the lower or upper endpoint is attained, \(M_T^{-}\) and \(M_T^{+}\) denote endpoint-attaining latent laws in \(\mathcal C_T(P_T,\rho)\), and \(z_T^{-},z_T^{+}\in K\) are their normalized projected endpoint coordinates.  The proposal's main question is whether \(\Theta_T(P_T,\rho)\) converges to \(\Theta_{\infty}(P_{\infty},\rho)\) in Hausdorff distance \(d_H\).

\begin{verbatim}
Locked parameters:
(qid, specialization, P, theta, identifying restrictions, bounding functional)
= (pid_reversible_dynamic_spillover, v1, P_infty, theta_infty(rho),
   R_rev,m,rho,Cesaro,
   [lower_theta_infty(P_infty,rho), upper_theta_infty(P_infty,rho)]).
P_infty is a projectively consistent infinite-panel observable law.
theta_infty(rho) is the Cesaro long-run sustained-treatment ATE.
R_rev,m,rho,Cesaro imposes consistency, bounded outcomes, finite memory,
positive recurrent reversible support, sustained-regime transition kernels,
and a global adjacent-spell carryover bound rho.
The bounding functional is the Hausdorff/Cesaro limit of finite-horizon sharp
intervals when the limit exists, otherwise the liminf/limsup interval pair.
\end{verbatim}

\section{Assumptions}
\begin{assumption}[Projective reversible panel law]\label{ass:projective}
The finite-dimensional marginals \(P_T\) of \(P_{\infty}\) are projectively consistent, \(D_t\) is not absorbing, \(\mathcal X\) is finite, and \(Y_t\in[0,1]\) for all \(t\).
\end{assumption}

\begin{assumption}[Finite memory and consistency]\label{ass:memory}
Potential outcomes are consistent with observed outcomes and depend on treatment paths only through the last \(m\) treatment statuses and the current \(X_t\).  Thus the sustained-regime rewards can be represented by the maps \(r_1,r_0:\mathcal X\to[0,1]\).
\end{assumption}

\begin{assumption}[Sustained-regime transition kernels]\label{ass:transition}
For \(a\in\{0,1\}\), the counterfactual transition kernel \(Q_a(x,x')\) of \(X_{t+1}\) under the sustained treatment \(a\) is identified from \(P_{\infty}\) on the reversible support and is time-homogeneous.
\end{assumption}

\begin{assumption}[Uniform geometric ergodicity]\label{ass:ergodic}
Each \(Q_a\) is irreducible and aperiodic on \(\mathcal X\), has stationary law \(\pi_a\), and satisfies \(\sup_x\|Q_a^k(x,\cdot)-\pi_a\|_{\mathrm{TV}}\le C\lambda^k\) for constants \(C<\infty\) and \(\lambda<1\) common to \(a=0,1\).
\end{assumption}

\begin{assumption}[Bounded adjacent-spell carryover]\label{ass:rho}
There is a scalar \(\rho\in[0,\infty]\) such that all latent reward maps compatible with \(P_{\infty}\) satisfy the global spillover inequalities used to construct the anchor intervals \(A_a(x;\rho)=[\ell_a(x;\rho),u_a(x;\rho)]\).  The inequalities compare only histories with the same current treatment status, so \(\rho=0\) means no residual carryover within a sustained arm, not no treatment effect.
\end{assumption}

\begin{assumption}[Nonempty dynamically compatible reward correspondence]\label{ass:compatible}
The set \(\mathcal G(P_{\infty},\rho)\) is nonempty, closed, and convex, and every \(r\in\mathcal G(P_{\infty},\rho)\) is induced by at least one latent infinite-panel law whose finite marginals agree with \(P_T\) and satisfy Assumptions~\ref{ass:projective}--\ref{ass:rho}.
\end{assumption}

\begin{assumption}[Robins endpoint positivity]\label{ass:robins-endpoint}
When \(\rho=0\), each sustained treatment history has the positivity and sequential exchangeability conditions needed for the Robins g-formula, and the anchor interval \(A_a(x;0)\) is the singleton containing the g-formula reward \(g_a(x)\).
\end{assumption}

\begin{assumption}[Unrestricted reward-cube endpoint]\label{ass:manski-endpoint}
When \(\rho=\infty\), the only remaining restriction on each sustained reward coordinate is \(r_a(x)\in[0,1]\), unless an observed sustained-history anchor imposes a singleton or interval restriction.
\end{assumption}

\begin{assumption}[Frontier-regular reversible sharpness class]\label{ass:face}
The law \(P_{\infty}\) belongs to the class \(\mathfrak P_{\mathrm{frontier}}(m,\rho)\) defined as follows.  The sequence \((\widehat{\mathcal Z}_T)_{T\ge1}\) consists of nonempty compact convex subsets of \(K\); every subsequential Hausdorff limit is the projection of at least one projectively consistent latent panel law satisfying Assumptions~\ref{ass:projective}--\ref{ass:rho}; and the stationary set \(\mathcal Z_\infty\) has nonempty exposed endpoint faces \(F_\infty^-\) and \(F_\infty^+\).  The class excludes degenerate single-horizon artifacts but does not impose time-homogeneity, stationarity, optimizer selection, endpoint convergence, or \(\Delta_T\to0\) on the latent reward faces.
\end{assumption}

\begin{assumption}[Two-face block-switching witness geometry]\label{ass:cycling}
For the diagnostic construction only, Assumption~\ref{ass:face} is replaced by the following finite-state geometry.  The state space is \(\mathcal X=\{0,1\}\), \(m=1\), and
\[
Q_0=Q_1=\begin{pmatrix}\varepsilon&1-\varepsilon\\ 1-\varepsilon&\varepsilon\end{pmatrix}
\]
for some \(\varepsilon\in(0,1/2)\), so the sustained kernels are irreducible, aperiodic, and geometrically ergodic with a negative second eigenvalue and stationary distribution \(\pi=(1/2,1/2)\).  The initial state law is \(\mu_0=(1,0)\).  The observed law has positive probability for both treatment statuses at both states and observed conditional mean outcomes in \((\delta,1-\delta)\) for a fixed \(\delta\in(0,1/6)\).  There are two nondegenerate projected reward faces \(F^A,F^B\subset K\) and constants \(\eta>0\), \(0<\alpha<1\), and increasing block endpoints \(1<T_1<T_2<\cdots\) with \(T_{k+1}/T_k\to\infty\) such that \(F^A\) contains
\[
z^A=(\pi,\pi,r_1^A,r_0^A),\qquad
r_1^A=(2\delta,2\delta),\quad r_0^A=(1-2\delta,1-2\delta),
\]
while \(F^B\) contains
\[
z^B=(\pi,\pi,r_1^B,r_0^B),\qquad
r_1^B=(1/2-\delta,1/2-\delta),\quad
r_0^B=(1/2+\delta,1/2+\delta).
\]
The lower support values on neighborhoods of these two faces differ by at least \(\eta\le 1-6\delta\).  The finite projected sharpness set satisfies
\[
d_H(\widehat{\mathcal Z}_{T_k},F^A)\le \alpha^{k}\quad\text{for odd }k,\qquad
d_H(\widehat{\mathcal Z}_{T_k},F^B)\le \alpha^{k}\quad\text{for even }k,
\]
and both faces are compatible with the same finite spillover radius \(\rho\).  The block endpoints are generated by a projectively consistent sequence of latent carryover coordinates that is observed only through \(P_T\)'s sharp compatibility inequalities, so endpoint switching is forced by the finite sharpness set itself rather than by choosing different optimizers on a fixed tied face.
\end{assumption}

\section{Main results}
\begin{proposition}[Background stationary reward support formula]\label{thm:stationary-reward}
Under Assumptions~\ref{ass:projective}--\ref{ass:compatible}, suppose further that the sustained-regime kernels \(Q_1,Q_0\) have stationary laws \(\pi_1,\pi_0\).  Then the stationary long-run support set is the interval
\[
\Theta_{\infty}(P_{\infty},\rho)
=\left[
\inf_{r\in\mathcal G(P_{\infty},\rho)}
  \{\pi_1 r_1-\pi_0 r_0\},
\sup_{r\in\mathcal G(P_{\infty},\rho)}
  \{\pi_1 r_1-\pi_0 r_0\}
\right],
\]
where \(\pi_a r_a=\sum_x\pi_a(x)r_a(x)\).  If \(\mathcal G(P_{\infty},\rho)\) is described by finitely many linear inequalities, both endpoints are attained at extreme points of \(\mathcal G(P_{\infty},\rho)\).
\end{proposition}

\paragraph{Why this is new and worth studying.}
(a) This proposition is not a novelty-bearing result; it is included as background notation for the Markov-reward support interval.  The closest published results are the average-reward and occupation-measure formulas in \citet{Puterman1994} and \citet{Altman1999}; the epsilon-gap for this proposal begins only after that formula, at the causal question of whether \(\mathcal G(P_\infty,\rho)\) is the sharp long-run projection of a projective reversible-spell latent-law problem.  (b) The concrete consumer is Conjecture~\ref{conj:cesaro-commutation}: without this stationary target, the finite-horizon intervals studied by \citet{Han2021} have no precise long-run object to converge to.

\begin{conjecture}[Conjecture 1: finite-window frontier for sharp reversible-spell bounds (NOT YET PROVED)]\label{conj:cesaro-commutation}
Under Assumptions~\ref{ass:projective}--\ref{ass:compatible}, \ref{ass:ergodic}, and \ref{ass:face}, the following statements are equivalent:
\[
\Psi_T=\Delta_T\to0
\]
and
\[
d_H\!\left(\Theta_T(P_T,\rho),\Theta_{\infty}(P_{\infty},\rho)\right)\to0.
\]
When these equivalent conditions hold, the stationary lift defect also vanishes, \(\Lambda_T\to0\), and
\[
\lim_{T\to\infty}\underline\theta_T
=\inf_{r\in\mathcal G(P_{\infty},\rho)}(\pi_1r_1-\pi_0r_0),
\qquad
\lim_{T\to\infty}\overline\theta_T
=\sup_{r\in\mathcal G(P_{\infty},\rho)}(\pi_1r_1-\pi_0r_0).
\]
Conversely, if \(\Delta_T\not\to0\), then on a nonempty open subclass of \(\mathfrak P_{\mathrm{frontier}}(m,\rho)\) sharing the same finite sustained-regime kernels and the same bounded-carryover radius, at least one endpoint sequence \((\underline\theta_T)\) or \((\overline\theta_T)\) fails to converge to the corresponding stationary support value.  If the finite-horizon lower and upper faces are represented by endpoint laws \(M_T^{-}\) and \(M_T^{+}\), then under \(\Delta_T\to0\) every limit point of the normalized projected endpoint coordinates \(z_T^{-}\) and \(z_T^{+}\) lies in \(F_{\infty}^{-}\) or \(F_{\infty}^{+}\), respectively, and admits an infinite-panel lift satisfying Assumption~\ref{ass:compatible}.
\end{conjecture}

\paragraph{Why this is new and worth studying.}
(a) The closest paper is \citet{Han2021}: it identifies finite dynamic sequence effects, but its theorem does not say when the sequence of sharp sets for reversible bounded-memory spells converges to a long-run ATE.  The epsilon-gap to \citet{ChernozhukovFernandezValHahnNewey2013} and \citet{JunLeeShin2016} is that their panel sharpness comes from time-homogeneity or stationarity imposed directly, while this conjecture gives a sharp if-and-only-if frontier for when finite reversible-spell causal compatibility polytopes generate the stationary lift.  The epsilon-gap to \citet{Schneider1993} and \citet{RockafellarWets1998} is that generic support-function stability does not define the observable endpoint diameter \(\Delta_{T,J}\), the causal stationary lift consequence \(\Lambda_T\to0\), or the open-subclass failure implication when \(\Delta_T\not\to0\).  (b) Resolving this conjecture would give analysts of \citet{DeChaisemartinDHaultfoeuille2024}-style nonabsorbing lagged-treatment panels a valid long-run sensitivity interval when they choose not to impose parallel trends.

\begin{conjecture}[Conjecture 2: block-switching sharpness failure (NOT YET PROVED)]\label{conj:face-cycling}
Assumptions~\ref{ass:projective}--\ref{ass:compatible} and \ref{ass:ergodic} do not by themselves imply Conjecture~\ref{conj:cesaro-commutation}.  There exists a finite-state reversible-spell law and a finite \(\rho\) satisfying Assumption~\ref{ass:cycling} such that \(Q_1\) and \(Q_0\) are uniformly geometrically ergodic, every finite \(\Theta_T(P_T,\rho)\) is sharp and nonempty, the observed law has nonconstant outcomes in both sustained arms, and the projected sharpness sets approach \(F^A\) along the odd block endpoints and \(F^B\) along the even block endpoints.  Consequently,
\[
\liminf_{T\to\infty}\underline\theta_T
<
\limsup_{T\to\infty}\underline\theta_T
\quad\text{or}\quad
\liminf_{T\to\infty}\overline\theta_T
<
\limsup_{T\to\infty}\overline\theta_T .
\]
In this example, the stationary reward interval from Proposition~\ref{thm:stationary-reward} is a valid subsequential envelope but is not the Hausdorff limit of the finite-horizon sharp identified sets.  Therefore the finite constrained-optimization workflow represented by \citet{KalouptsidiKitamuraLimaSouzaRodrigues2020} and the automated finite sharp-bounds workflow of \citet{DuarteFinkelsteinKnoxMummoloShpitser2024} would be invalid if their horizon-by-horizon endpoints were summarized by a single stationary interval without first reporting the frontier diagnostic \(\Delta_T\).
\end{conjecture}

\paragraph{Why this is new and worth studying.}
(a) The closest published long-run dynamic-treatment papers, \citet{Han2021} and \citet{HeckmanHumphriesVeramendi2016}, make finite dynamic timing and continuation values explicit but do not give a projected-face diagnostic for when horizon-indexed sharp sets define a long-run object.  The closest mathematical stability references, \citet{Schneider1993} and \citet{RockafellarWets1998}, explain why a fixed compact-convex limit would control support values; this conjecture supplies the causal epsilon-gap, namely a projectively consistent reversible-spell sharpness construction in which the finite causal feasible sets themselves switch between two endpoint faces.  The closest finite sharp-bounds papers, \citet{DuarteFinkelsteinKnoxMummoloShpitser2024} and \citet{KalouptsidiKitamuraLimaSouzaRodrigues2020}, solve finite optimization problems; this conjecture names the exact extrapolation those workflows must avoid, namely replacing a horizon-by-horizon endpoint sequence by one stationary interval when \(\Delta_T\) has no tail limit.  (b) The concrete consumer is \citet{DiemerShiSwanson2024}: the diagnostic specifies when sustained dynamic partial-ID bounds require a liminf/limsup report rather than a single long-run interval.

\section{Sanity-check corollaries}
\begin{corollary}[Zero-carryover Robins reduction]\label{cor:robins}
Under Assumptions~\ref{ass:projective}--\ref{ass:face} and \ref{ass:robins-endpoint}, if \(\rho=0\), then
\[
\Theta_{\infty}(P_{\infty},0)
=\left\{\sum_x\pi_1(x)g_1(x)-\sum_x\pi_0(x)g_0(x)\right\}.
\]
Thus Conjecture~\ref{conj:cesaro-commutation} collapses to the long-run Robins g-formula contrast rather than to a zero treatment effect.
\end{corollary}

\begin{corollary}[Unrestricted dynamic Manski endpoint]\label{cor:manski}
Under Assumptions~\ref{ass:projective}--\ref{ass:compatible} and \ref{ass:manski-endpoint}, if \(\rho=\infty\) and no observed sustained-history anchor tightens the reward cube, then \(\Theta_{\infty}(P_{\infty},\infty)=[-1,1]\) for the sustained-treatment contrast.  With observed sustained-history anchors, the endpoint tightens only through those anchor singleton or interval restrictions, matching the Diemer--Manski unrestricted logic.
\end{corollary}

\section{Proof-sketch route}
Proposition~\ref{thm:stationary-reward} reduces to the finite-state ergodic theorem for Markov reward processes plus elementary compact-convex support-function algebra.  For Conjecture~\ref{conj:cesaro-commutation}, the proposed route is to write each finite endpoint as the support value of \(\widehat{\mathcal Z}_T\); one direction is the Cauchy criterion for the two scalar endpoint sequences plus the definition of the stationary lift defect \(\Lambda_T\).  The reverse direction is the substantive step: on \(\mathfrak P_{\mathrm{frontier}}(m,\rho)\), endpoint convergence to \(\Theta_\infty\) must force vanishing endpoint diameter and rule out a nonzero stationary lift defect by separating any nonliftable subsequential face with a supporting hyperplane.  The causal work is this separation/liftability argument, because limiting projected faces must be shown to come from infinite-panel latent laws satisfying the same bounded-carryover restrictions, not merely from abstract compact-set convergence.  Uniform ergodicity supplies the boundary-term bound once the stationary lift is identified, and Assumption~\ref{ass:compatible} supplies the infinite-panel lift of limiting faces.  The boundary cases needing explicit algebra are \(m=1\), \(|\mathcal X|=2\), and \(T=3,4\), where one can enumerate \(\widehat{\mathcal Z}_T\), \(F_T^s\), \(\Delta_{T,J}\), and \(\Gamma_T^s(\varepsilon)\).  Conjecture~\ref{conj:face-cycling} should be attacked by constructing the block sequence in Assumption~\ref{ass:cycling}: verify the displayed rewards \(z^A,z^B\) keep all observed conditional means in \((\delta,1-\delta)\), compute the two lower support values \(4\delta-1\) and \(-2\delta\), and show that the sharp finite compatibility inequalities force odd block endpoints within \(\alpha^k\) of \(F^A\) and even block endpoints within \(\alpha^k\) of \(F^B\).

\section{Honest scope flags}
Proposition~\ref{thm:stationary-reward} is background finite-state Markov-reward support algebra once the stationary reward correspondence is fixed.  Conjecture~\ref{conj:cesaro-commutation} is open and now asserts an if-and-only-if frontier only on the named class \(\mathfrak P_{\mathrm{frontier}}(m,\rho)\); outside that class, the proposal does not claim endpoint convergence forces full Hausdorff convergence of all projected polytopes.  Conjecture~\ref{conj:face-cycling} is also open and is now scoped as a named-warning result for finite-horizon sharp-bounds workflows, with endpoint separation forced by horizon-indexed causal compatibility sets rather than by opposite signs of a decaying boundary correction.  This pivot intentionally withdraws the stronger dead-angle claim that a universal closed-form nested g-formula bridge exists for every \(\rho\); the present proposal asks the prior question of when such finite-horizon intervals have a legitimate long-run limit at all.

\section{Iteration trail}
\begin{itemize}[leftmargin=*]
\item v1 (pivot) -- initial draft for angle 2, replacing the rejected finite-radius sharpness-witness kernel with a Cesaro-limit and face-cycling kernel anchored to \citet{Han2021}; prior verdict: none.
\item v2 (revise) -- fixed the Han--Yang DOI, replaced informal face stability by a projected occupation-reward face condition with liftability, and recast face-cycling as a diagnostic rather than a strawman refutation; prior verdict: REVISE.
\item v3 (revise) -- promoted Conjecture~\ref{conj:cesaro-commutation} to an observable necessary-and-sufficient Hausdorff-face stability criterion, corrected the Han--Yang comparison, and separated the two-face witness assumptions from the claimed endpoint subsequence failure; prior verdict: REVISE.
\item v4 (revise) -- demoted the standard Markov reward formula to background, corrected the de Chaisemartin--D'Haultfoeuille citation metadata, added convex-analysis prior art, replaced the false endpoint-to-face equivalence by a tail-stability sufficiency claim, and strengthened the failure witness to block-switching finite sharpness sets; prior verdict: REVISE.
\item v5 (revise) -- promoted Conjecture~\ref{conj:cesaro-commutation} from sufficiency to a finite-window if-and-only-if frontier, named \citet{KalouptsidiKitamuraLimaSouzaRodrigues2020} and \citet{DuarteFinkelsteinKnoxMummoloShpitser2024} as finite sharp-bounds workflows whose stationary extrapolation is invalidated by Conjecture~\ref{conj:face-cycling}, and added the direct CFHN time-homogeneity comparison; prior verdict: REVISE.
\end{itemize}

\section*{Bibliography}
\begin{thebibliography}{99}
\bibitem[Chernozhukov et~al.(2013)]{ChernozhukovFernandezValHahnNewey2013}
Chernozhukov, V., Fernandez-Val, I., Hahn, J., and Newey, W. (2013).
\newblock Average and quantile effects in nonseparable panel models.
\newblock \emph{Econometrica} 81(2), 535--580.
\newblock doi:10.3982/ECTA8405.

\bibitem[de Chaisemartin and D'Haultfoeuille(2024)]{DeChaisemartinDHaultfoeuille2024}
de Chaisemartin, C. and D'Haultfoeuille, X. (2024).
\newblock Difference-in-differences estimators of intertemporal treatment effects.
\newblock \emph{Review of Economics and Statistics}, February 9, 2024, 1--45.
\newblock doi:10.1162/rest\_a\_01414.

\bibitem[Diemer et~al.(2024)]{DiemerShiSwanson2024}
Diemer, E. W., Shi, J., and Swanson, S. A. (2024).
\newblock Partial identification of the effects of sustained treatment strategies.
\newblock \emph{Epidemiology} 35(3).
\newblock doi:10.1097/EDE.0000000000001721.

\bibitem[Duarte et~al.(2024)]{DuarteFinkelsteinKnoxMummoloShpitser2024}
Duarte, G., Finkelstein, N., Knox, D., Mummolo, J., and Shpitser, I. (2024).
\newblock An automated approach to causal inference in discrete settings.
\newblock \emph{Journal of the American Statistical Association} 119(547), 1778--1793.
\newblock doi:10.1080/01621459.2023.2216909.

\bibitem[Han(2021)]{Han2021}
Han, S. (2021).
\newblock Identification in nonparametric models for dynamic treatment effects.
\newblock \emph{Journal of Econometrics} 225(2), 132--147.
\newblock doi:10.1016/j.jeconom.2019.08.014.

\bibitem[Han and Yang(2024)]{HanYang2024}
Han, S. and Yang, S. (2024).
\newblock A computational approach to identification of treatment effects for policy evaluation.
\newblock \emph{Journal of Econometrics} 240(1), 105680.
\newblock doi:10.1016/j.jeconom.2024.105680.

\bibitem[Heckman et~al.(2016)]{HeckmanHumphriesVeramendi2016}
Heckman, J. J., Humphries, J. E., and Veramendi, G. (2016).
\newblock Dynamic treatment effects.
\newblock \emph{Journal of Econometrics} 191(2), 276--292.
\newblock doi:10.1016/j.jeconom.2015.12.001.

\bibitem[Jun, Lee, and Shin(2016)]{JunLeeShin2016}
Jun, S. J., Lee, Y., and Shin, Y. (2016).
\newblock Treatment effects with unobserved heterogeneity: A set identification approach.
\newblock \emph{Journal of Business and Economic Statistics} 34(2), 302--311.
\newblock doi:10.1080/07350015.2015.1044008.

\bibitem[Kalouptsidi et~al.(2020)]{KalouptsidiKitamuraLimaSouzaRodrigues2020}
Kalouptsidi, M., Kitamura, Y., Lima, L., Souza-Rodrigues, E., and others (2020).
\newblock Partial identification and inference for dynamic models and counterfactuals.
\newblock \emph{NBER Working Paper} 26761.
\newblock doi:10.3386/w26761.

\bibitem[Altman(1999)]{Altman1999}
Altman, E. (1999).
\newblock \emph{Constrained Markov Decision Processes}.
\newblock Chapman and Hall/CRC.

\bibitem[Puterman(1994)]{Puterman1994}
Puterman, M. L. (1994).
\newblock \emph{Markov Decision Processes: Discrete Stochastic Dynamic Programming}.
\newblock Wiley.

\bibitem[Rockafellar and Wets(1998)]{RockafellarWets1998}
Rockafellar, R. T. and Wets, R. J.-B. (1998).
\newblock \emph{Variational Analysis}.
\newblock Springer.

\bibitem[Robins(1986)]{Robins1986}
Robins, J. M. (1986).
\newblock A new approach to causal inference in mortality studies with a sustained exposure period -- application to control of the healthy worker survivor effect.
\newblock \emph{Mathematical Modelling} 7(9--12), 1393--1512.
\newblock doi:10.1016/0270-0255(86)90088-6.

\bibitem[Schneider(1993)]{Schneider1993}
Schneider, R. (1993).
\newblock \emph{Convex Bodies: The Brunn--Minkowski Theory}.
\newblock Cambridge University Press.
\end{thebibliography}

\end{document}
